\documentclass[12pt]{article}

\usepackage[utf8]{inputenc}
\usepackage[margin=1in]{geometry}
\usepackage{setspace}
\doublespacing

\title{A Simple Recipe for Belonging: A Film Adaptation Pitch}

\author{%
  Ma Toan Bach\\
  LSO100KCC\\
  Film Adaptation Pitch
}

\date{}

\begin{document}

\maketitle

My film adaptation of Madeleine Thien's ``Simple Recipes'' would be an intimate contemporary drama set in Vancouver. The story follows a daughter remembering her Malaysian-Chinese immigrant father, especially the way he cooks rice, prepares fish, and fills the family kitchen with ritual, skill, and authority. At first, the father appears almost magical to the child narrator. He knows how to measure rice by touch, carve fruit, cook without hesitation, and turn ordinary meals into moments of family warmth. However, the memory darkens when the brother rejects the fish dinner, insults his father with racist language, and is violently beaten. The film would preserve the story's central themes: immigration, cultural inheritance, language loss, family love, shame, and the painful difficulty of belonging to more than one world.

I would update the story to present-day Vancouver while keeping its emotional structure and domestic focus. The film would move between the adult narrator's clean, quiet apartment and her childhood kitchen, using memory as the main storytelling device. The adult narrator would rarely explain events directly; instead, the camera would return to sensory details: cloudy rice water, steam on glasses, oil crackling in a wok, a fish moving weakly in the sink, and the sound of a fork scraping against a plate. The tone would begin warm and almost nostalgic, then slowly become tense. This style would show how food can hold both comfort and trauma. A simple recipe becomes a way of remembering love, but also a reminder that family history cannot be cleaned as easily as rice.

The casting should avoid celebrity distraction and focus on actors who can communicate emotion through silence. The father must be portrayed as loving, skilled, proud, and frightening, because the power of the story depends on the narrator's inability to separate these parts of him. He is not a simple villain. He is an immigrant father who carries displacement, pride, and wounded authority into the home. The mother would be played as a tired but emotionally alert figure, someone who works outside the house and then returns to manage the family's pain. The brother would represent a different response to immigration: rejection, anger, and the desire not to be marked as different. The daughter would be shown first as admiring and dependent, then as slowly aware that love can become mixed with guilt and grief.

The key scene that must appear in the film is the fish dinner. I would film it mostly through close-ups and sound. The father would prepare the fish with careful hands, while the brother sits at the table dirty from soccer, already tense and separate from the family ritual. The camera would linger on the fish's eye, the steam, the chopsticks, and the brother's discomfort. When the father offers him fish and jokes, ``Take a wok on the wild side,'' the line should feel like an attempt at connection. The brother's refusal and racist outburst would then break the scene open. I would avoid melodramatic music and let the silence afterward feel unbearable, especially as the violence moves from the dinner table to the bedroom.

This adaptation would matter to contemporary Canadian audiences because it shows multiculturalism as something lived privately, not just celebrated publicly. Canada often presents diversity through food, festivals, and inclusion, but Thien's story asks what happens inside families when language is lost, children assimilate, and parents feel rejected by the country their children now call home. The brother's insult is devastating because it turns public racism into family speech. The narrator's memory shows that belonging is not simple: it can involve love for parents, shame about difference, fear of violence, and grief for what cannot be repaired. My film would preserve that complexity by making the kitchen both a place of nourishment and a place where cultural conflict becomes impossible to ignore.

\end{document}
